\documentclass[10pt,a4paper]{article}
\usepackage[margin=1.6cm]{geometry}
% common.tex - shared preamble for all ELL7282 A2 documents
\usepackage[T1]{fontenc}
\usepackage{lmodern}
\usepackage{amsmath,amssymb}
\usepackage{booktabs,longtable,tabularx,array,multirow}
\usepackage[table]{xcolor}
\usepackage{bytefield}
\usepackage{tikz}
\usetikzlibrary{arrows.meta,positioning,fit,calc,shapes.geometric}
\usepackage{pdflscape}
\usepackage{enumitem}
\usepackage{caption}
\usepackage[hidelinks]{hyperref}
\setlist{nosep,leftmargin=1.4em}
\setlength{\parskip}{3pt}
\setlength{\parindent}{0pt}

\newcommand{\project}{WM-8 five-stage pipelined washing-machine microcontroller}
\newcommand{\course}{ELL7282/ELL782 Introduction to Computer Architecture -- Assignment 2}
\newcommand{\pass}{\textcolor{green!45!black}{\textbf{PASS}}}
\newcommand{\fail}{\textcolor{red!75!black}{\textbf{FAIL}}}
% invalid (V=0) pipeline entry: grey with an empty-set marker
\newcommand{\bub}[1]{\textcolor{black!45}{$\varnothing$\,\texttt{#1}}}
\newcommand{\sig}[1]{\texttt{#1}}
\newcommand{\hex}[1]{\texttt{0x#1}}

% cycle-trace table (used by generated/trace*.tex)
\newenvironment{tracetab}{%
  \scriptsize\setlength{\tabcolsep}{2.5pt}\renewcommand{\arraystretch}{1.15}%
  \begin{longtable}{r p{3.0cm} p{3.0cm} p{3.0cm} p{3.0cm} p{3.0cm} c c c c p{4.6cm}}
  \toprule
  Cyc & IF (fetch PC) & ID (IF/ID) & EX (ID/EX) & MEM (EX/MEM) & WB (MEM/WB) &
  Stall & Flush & Fwd A & Fwd B & Architectural writes (commit at end of cycle)\\
  \midrule\endhead
  \bottomrule\endfoot}{%
  \end{longtable}}

% validation table (used by generated/val_g*.tex)
\newenvironment{valtab}{%
  \footnotesize\renewcommand{\arraystretch}{1.2}%
  \begin{longtable}{|p{0.6cm}|p{3.7cm}|p{4.2cm}|p{5.6cm}|p{1.0cm}|}
  \hline
  \rowcolor{black!8}\textbf{ID} & \textbf{Initial inputs / state, stimulus} & \textbf{Expected result} &
  \textbf{Observed result} & \textbf{P/F}\\ \hline\endhead}{%
  \end{longtable}}

% program listing (used by generated/listing_*.tex)
\newenvironment{listtab}{%
  \footnotesize\setlength{\tabcolsep}{4pt}\renewcommand{\arraystretch}{1.05}%
  \begin{longtable}{l l c l l p{6.2cm}}
  \toprule
  Addr & Hex & Fmt & Label & Assembly & Comment\\
  \midrule\endhead
  \bottomrule\endfoot}{%
  \end{longtable}}

\begin{document}
\begin{center}
{\Large\textbf{Cycle-by-Cycle Pipeline Traces (Sections 8.1--8.4)}}\\[2pt]
\project\ -- \course
\end{center}

\section{How the traces were produced}
The test benches \texttt{tests/t1\_hazards.circ} (Sections 8.1--8.3) and \texttt{tests/t2\_branch\_halt.circ} (Section 8.4) contain
the unmodified \texttt{cpu} and \texttt{mmio} sub-circuits, the instruction ROM holding the test program, the data RAM, a
power-on reset (first cycle), a cycle counter and the pipeline-monitor outputs of the CPU. These outputs are the PC of each
pipeline register and its valid bit, \texttt{STALL}, \texttt{FLUSH}, \texttt{HALTED}, \texttt{FWD\_A}/\texttt{FWD\_B}, the
write-back port (\texttt{WB\_WE}, \texttt{WB\_RD}, \texttt{WB\_DATA}) and the data-bus write (\texttt{D\_WE},
\texttt{D\_ADDR}, \texttt{D\_WDATA}). One row is one clock cycle: the state after rising edge \emph{n}. Each instruction cell
shows the instruction address and source text, taken from the \texttt{PC} field travelling with it. The
rows were recorded by cycle-accurate simulation of the test-bench \texttt{.circ} netlists (\texttt{tools/circsim.py}); the
raw rows are in \texttt{tests/expected/*.txt}, in the same tab-separated format that Logisim prints with
\texttt{java -jar logisim-generic-2.7.1.jar tests/t1\_hazards.circ -tty table}, so they can be compared directly.

\textbf{Notation.} \texttt{08 LD R1,[R2]} = valid entry; \bub{09 ADD R3,R1} = entry with \texttt{V}=0 (bubble or squashed
instruction; its stale fields are shown only for orientation). \emph{Stall} = \texttt{STALL} (PC and IF/ID hold at the end of the
cycle, ID/EX receives a bubble). \emph{Flush} = taken branch/jump or HALT in EX (IF/ID and ID/EX invalidated at the end of the
cycle). \emph{Fwd A/B}: source of the EX operand \texttt{ra}/\texttt{rb} (\texttt{--} = value read in ID).
\emph{Architectural writes}: register write-back and data-bus stores that commit at the rising edge \emph{ending} the cycle.

\section{Section 8.1 -- ALU forwarding without stall}
Program: \texttt{ADD R1,R2 (06); SUB R3,R1 (07)} with R1=01, R2=10, R3=03. In \textbf{cycle 10} of the table below,
\texttt{SUB R3,R1} is in EX while \texttt{ADD R1,R2} is in MEM. \texttt{FWD\_B}=EX/MEM supplies R1=\hex{11}, there is no stall,
and R3$\leftarrow$\hex{F2} is written back in cycle 12.

\begin{landscape}
\section{Section 8.2 -- load-use dependence (exactly one stall)}
Program: \texttt{LD R1,[R2] (08); ADD R3,R1 (09)}, RAM[\hex{10}]=\hex{03} written earlier by \texttt{ST R3,[R2]}.
\begin{tracetab}
8 & \texttt{07 SUB R3, R1} & \texttt{06 ADD R1, R2} & \texttt{05 NOP} & \texttt{04 ST R3, [R2]} & \texttt{03 LDI R4, 0x20} & 0 & 0 & -- & -- & R4$\leftarrow$20; RAM[10]$\leftarrow$03 \\
9 & \texttt{08 LD R1, [R2]} & \texttt{07 SUB R3, R1} & \texttt{06 ADD R1, R2} & \texttt{05 NOP} & \texttt{04 ST R3, [R2]} & 0 & 0 & -- & -- & -- \\
10 & \texttt{09 ADD R3, R1} & \texttt{08 LD R1, [R2]} & \texttt{07 SUB R3, R1} & \texttt{06 ADD R1, R2} & \texttt{05 NOP} & 0 & 0 & -- & EX/MEM & -- \\
11 & \texttt{0A ADD R1, R2} & \texttt{09 ADD R3, R1} & \texttt{08 LD R1, [R2]} & \texttt{07 SUB R3, R1} & \texttt{06 ADD R1, R2} & \textbf{1} & 0 & -- & -- & R1$\leftarrow$11 \\
12 & \texttt{0A ADD R1, R2} & \texttt{09 ADD R3, R1} & \bub{09 ADD R3, R1} & \texttt{08 LD R1, [R2]} & \texttt{07 SUB R3, R1} & 0 & 0 & -- & -- & R3$\leftarrow$F2 \\
13 & \texttt{0B SUB R1, R3} & \texttt{0A ADD R1, R2} & \texttt{09 ADD R3, R1} & \bub{09 ADD R3, R1} & \texttt{08 LD R1, [R2]} & 0 & 0 & -- & MEM/WB & R1$\leftarrow$03 \\
14 & \texttt{0C ST R1, [R4]} & \texttt{0B SUB R1, R3} & \texttt{0A ADD R1, R2} & \texttt{09 ADD R3, R1} & \bub{09 ADD R3, R1} & 0 & 0 & -- & -- & -- \\
15 & \texttt{0D LD R5, [R4]} & \texttt{0C ST R1, [R4]} & \texttt{0B SUB R1, R3} & \texttt{0A ADD R1, R2} & \texttt{09 ADD R3, R1} & 0 & 0 & EX/MEM & MEM/WB & R3$\leftarrow$F5 \\

\end{tracetab}
\vspace{-4mm}
\begin{itemize}
\item \textbf{Cycle 11}: \texttt{LD R1} is in EX and \texttt{ADD R3,R1} (which reads R1 via \texttt{rb}, \texttt{UB}=1) is in ID, so
  \texttt{STALL}=1. This is the only stall cycle.
\item \textbf{Cycle 12}: the PC (\hex{0A}) and IF/ID (\texttt{ADD R3,R1}) are frozen and ID/EX holds a bubble
  (\bub{09}). Older instructions advance: \texttt{LD} reads RAM in MEM and \texttt{SUB} writes R3=\hex{F2} in WB. The stalled
  \texttt{ADD} reads that R3 through the WB$\rightarrow$ID bypass.
\item \textbf{Cycle 13}: \texttt{ADD R3,R1} is in EX and \texttt{LD} is in WB, so \texttt{FWD\_B}=MEM/WB supplies R1=\hex{03}. The load
  was not forwarded from EX/MEM in cycle 12, when it was not yet ready. \textbf{Cycle 15}: R3$\leftarrow$\hex{F2}+\hex{03}=\hex{F5}.
\end{itemize}

\section{Section 8.3 -- youngest producer forwarded to store data}
Program: \texttt{ADD R1,R2 (0A); SUB R1,R3 (0B); ST R1,[R4] (0C)} with R1=03, R2=10, R3=F5, R4=20.
\begin{tracetab}
12 & \texttt{0A ADD R1, R2} & \texttt{09 ADD R3, R1} & \bub{09 ADD R3, R1} & \texttt{08 LD R1, [R2]} & \texttt{07 SUB R3, R1} & 0 & 0 & -- & -- & R3$\leftarrow$F2 \\
13 & \texttt{0B SUB R1, R3} & \texttt{0A ADD R1, R2} & \texttt{09 ADD R3, R1} & \bub{09 ADD R3, R1} & \texttt{08 LD R1, [R2]} & 0 & 0 & -- & MEM/WB & R1$\leftarrow$03 \\
14 & \texttt{0C ST R1, [R4]} & \texttt{0B SUB R1, R3} & \texttt{0A ADD R1, R2} & \texttt{09 ADD R3, R1} & \bub{09 ADD R3, R1} & 0 & 0 & -- & -- & -- \\
15 & \texttt{0D LD R5, [R4]} & \texttt{0C ST R1, [R4]} & \texttt{0B SUB R1, R3} & \texttt{0A ADD R1, R2} & \texttt{09 ADD R3, R1} & 0 & 0 & EX/MEM & MEM/WB & R3$\leftarrow$F5 \\
16 & \texttt{0E LDI R7, 0xF5} & \texttt{0D LD R5, [R4]} & \texttt{0C ST R1, [R4]} & \texttt{0B SUB R1, R3} & \texttt{0A ADD R1, R2} & 0 & 0 & EX/MEM & -- & R1$\leftarrow$13 \\
17 & \texttt{0F ST R5, [R7]} & \texttt{0E LDI R7, 0xF5} & \texttt{0D LD R5, [R4]} & \texttt{0C ST R1, [R4]} & \texttt{0B SUB R1, R3} & 0 & 0 & -- & -- & R1$\leftarrow$1E; RAM[20]$\leftarrow$1E \\
18 & \texttt{10 HALT} & \texttt{0F ST R5, [R7]} & \texttt{0E LDI R7, 0xF5} & \texttt{0D LD R5, [R4]} & \texttt{0C ST R1, [R4]} & 0 & 0 & -- & -- & -- \\

\end{tracetab}
\vspace{-4mm}
\begin{itemize}
\item \textbf{Cycle 15}: \texttt{SUB R1,R3} uses both forwarding paths at once: R1 from EX/MEM (\texttt{ADD}, \hex{13}) and R3
  from MEM/WB (\texttt{ADD R3,R1}, \hex{F5}).
\item \textbf{Cycle 16}: \texttt{ST R1,[R4]} is in EX. \emph{Two} in-flight producers of R1 exist: \texttt{SUB} in EX/MEM
  (\hex{1E}) and the older \texttt{ADD} in MEM/WB (\hex{13}). \texttt{FWD\_A}=EX/MEM, so the youngest wins and the store data is
  \hex{1E}. The address R4 comes from the register file (\texttt{FWD\_B}=\texttt{--}).
\item \textbf{Cycle 17}: the store commits once in MEM, RAM[\hex{20}]$\leftarrow$\hex{1E} (not the stale \hex{13}). The next
  instruction \texttt{LD R5,[R4]} reads back \hex{1E}, which the program then copies to \texttt{DISPLAY}.
\end{itemize}

\section{Section 8.4 -- taken branch squashes a wrong-path actuator store}
Program: \texttt{SUB R1,R1 (05); BEQZ R1,target1 (06); ST R2,[R4] (07); ADD R3,R4 (08); target1: (09)}
with R2=\hex{21} (would open the valve and lock the door) and R4=\hex{F4} (\texttt{Rvalve} = actuator latch).
\begin{tracetab}
5 & \texttt{04 LDI R7, 0xF5} & \texttt{03 LDI R1, 0x09} & \texttt{02 LDI R3, 0x07} & \texttt{01 LDI R4, 0xF4} & \texttt{00 LDI R2, 0x21} & 0 & 0 & -- & -- & R2$\leftarrow$21 \\
6 & \texttt{05 SUB R1, R1} & \texttt{04 LDI R7, 0xF5} & \texttt{03 LDI R1, 0x09} & \texttt{02 LDI R3, 0x07} & \texttt{01 LDI R4, 0xF4} & 0 & 0 & -- & -- & R4$\leftarrow$F4 \\
7 & \texttt{06 BEQZ R1, target1} & \texttt{05 SUB R1, R1} & \texttt{04 LDI R7, 0xF5} & \texttt{03 LDI R1, 0x09} & \texttt{02 LDI R3, 0x07} & 0 & 0 & -- & -- & R3$\leftarrow$07 \\
8 & \texttt{07 ST R2, [R4]} & \texttt{06 BEQZ R1, target1} & \texttt{05 SUB R1, R1} & \texttt{04 LDI R7, 0xF5} & \texttt{03 LDI R1, 0x09} & 0 & 0 & MEM/WB & MEM/WB & R1$\leftarrow$09 \\
9 & \texttt{08 ADD R3, R4} & \texttt{07 ST R2, [R4]} & \texttt{06 BEQZ R1, target1} & \texttt{05 SUB R1, R1} & \texttt{04 LDI R7, 0xF5} & 0 & \textbf{1} & EX/MEM & -- & R7$\leftarrow$F5 \\
10 & \texttt{09 BNEZ R1, bad} & \bub{08 ADD R3, R4} & \bub{07 ST R2, [R4]} & \texttt{06 BEQZ R1, target1} & \texttt{05 SUB R1, R1} & 0 & 0 & -- & -- & R1$\leftarrow$00 \\
11 & \texttt{0A BNEZ R3, target2} & \texttt{09 BNEZ R1, bad} & \bub{08 ADD R3, R4} & \bub{07 ST R2, [R4]} & \texttt{06 BEQZ R1, target1} & 0 & 0 & -- & -- & -- \\
12 & \texttt{0B HALT} & \texttt{0A BNEZ R3, target2} & \texttt{09 BNEZ R1, bad} & \bub{08 ADD R3, R4} & \bub{07 ST R2, [R4]} & 0 & 0 & -- & -- & -- \\
13 & \texttt{0C ST R2, [R4]} & \texttt{0B HALT} & \texttt{0A BNEZ R3, target2} & \texttt{09 BNEZ R1, bad} & \bub{08 ADD R3, R4} & 0 & \textbf{1} & -- & -- & -- \\
14 & \texttt{0D ST R3, [R7]} & \bub{0C ST R2, [R4]} & \bub{0B HALT} & \texttt{0A BNEZ R3, target2} & \texttt{09 BNEZ R1, bad} & 0 & 0 & -- & -- & -- \\

\end{tracetab}
\vspace{-4mm}
\begin{itemize}
\item \textbf{Cycle 9}: \texttt{BEQZ R1} is in EX and tests R1, forwarded from EX/MEM (\texttt{SUB R1,R1} = 0). The branch is
  taken, so \texttt{FLUSH}=1 and \texttt{PC}$\leftarrow$\hex{09}. At the same edge the wrong-path \texttt{ST R2,[R4]} (in ID)
  and \texttt{ADD R3,R4} (in IF) are written into ID/EX and IF/ID with \texttt{V}=0.
\item \textbf{Cycles 10--12}: the squashed entries keep moving down the pipe as bubbles. \texttt{ST R2,[R4]} reaches MEM in
  cycle 11 with \texttt{V}=0, so \texttt{D\_WE}=0, the actuator latch stays \hex{00} and \texttt{VALVE\_IN} never rises.
  \texttt{ADD R3,R4} never writes R3. Branch penalty = 2 cycles (09 is fetched in cycle 10).
\item \textbf{Cycle 12}: \texttt{BNEZ R1,bad} is not taken (no flush). \textbf{Cycle 13}: \texttt{BNEZ R3,target2} is taken
  and squashes the wrong-path \texttt{HALT} at \hex{0B}, which therefore never stops the machine (see full trace, Appendix B).
\end{itemize}

\end{landscape}

\section{Source-use detection (no stalls on fields that are not read)}
The decoder emits \texttt{UA} (instruction reads \texttt{ra}=\texttt{IR[11:9]}) and \texttt{UB} (reads \texttt{rb}=\texttt{IR[8:6]}):
\begin{center}\small
\begin{tabular}{l l l}
\toprule
\texttt{UA}=1 & \texttt{ADD, SUB, AND, SGEU, ADDI, ST} (data), \texttt{BEQZ, BNEZ} (tested register) & \\
\texttt{UB}=1 & \texttt{MOV, ADD, SUB, AND, SGEU, LD} (address), \texttt{ST} (address) & \\
neither & \texttt{NOP, LDI, JMP, HALT}, unassigned \hex{D}/\hex{E} & \\
\bottomrule
\end{tabular}
\end{center}
The load-use check is \texttt{IFID.V}$\wedge$\texttt{IDEX.V}$\wedge$\texttt{IDEX.MR}$\wedge$[(\texttt{UA}$\wedge$\texttt{RA}=\texttt{IDEX.RA})
$\vee$(\texttt{UB}$\wedge$\texttt{RB}=\texttt{IDEX.RA})]. Without \texttt{UA}/\texttt{UB}, the following cases would stall although
they read no register: \texttt{LDI R1,..} after \texttt{LD R1,..} (\texttt{ra} field equal), \texttt{JMP 0x4C} (bits 8:6 of the
target = 001), and \texttt{ADDI R3,0x40} or \texttt{BEQZ R2,0x47} (imm bits 7:6 and reserved bit 8 form a fake \texttt{rb}).
\texttt{MOV Rd,Rs} and \texttt{LD Rd,[Rs]} do not read \texttt{Rd}, so they must not stall on it. Forwarding selects are
qualified the same way (\texttt{USE\_A}=\texttt{IDEX.V}$\wedge$\texttt{UA}), so the monitor never reports spurious forwarding for
bubbles or unused fields. Bubbles (\texttt{IFID.V}=0 or \texttt{IDEX.V}=0) never cause a stall.

\section{A store commits at most once}
The data-bus write strobe is \texttt{D\_WE}=\texttt{EXMEM.V}$\wedge$\texttt{EXMEM.MW}, and the RAM, \texttt{ACT},
\texttt{DISPLAY}, \texttt{T\_LIMIT} and \texttt{T\_CTRL} registers sample it only at the rising edge that ends the MEM cycle.
Stalls freeze only \texttt{PC} and IF/ID. EX/MEM and MEM/WB load new contents every cycle, so a store occupies EX/MEM for
exactly one cycle and the stalled copy in ID is not yet in MEM. On a flush or halt, the squashed copies are written with
\texttt{V}=0. In every trace each store appears in the \emph{writes} column once (e.g.\ RAM[\hex{20}] in cycle 17
above), and squashed stores never appear (cycle 11 of 8.4).

\appendix
\begin{landscape}
\section{Full trace of \texttt{t1\_hazards} (cycles 0--24)}
Cycles 0--1: power-on reset, then the first fetch. Cycle 20: \texttt{HALT} in EX flushes the younger \texttt{LDI R6} and freezes the PC.
From cycle 21 \texttt{HALTED}=1, and older instructions drain.
\begin{tracetab}
0 & \texttt{00 LDI R1, 0x01} & \bub{00 LDI R1, 0x01} & \bub{00 LDI R1, 0x01} & \bub{00 LDI R1, 0x01} & \bub{00 LDI R1, 0x01} & 0 & 0 & -- & -- & -- \\
1 & \texttt{00 LDI R1, 0x01} & \bub{00 LDI R1, 0x01} & \bub{00 LDI R1, 0x01} & \bub{00 LDI R1, 0x01} & \bub{00 LDI R1, 0x01} & 0 & 0 & -- & -- & -- \\
2 & \texttt{01 LDI R2, 0x10} & \texttt{00 LDI R1, 0x01} & \bub{00 LDI R1, 0x01} & \bub{00 LDI R1, 0x01} & \bub{00 LDI R1, 0x01} & 0 & 0 & -- & -- & -- \\
3 & \texttt{02 LDI R3, 0x03} & \texttt{01 LDI R2, 0x10} & \texttt{00 LDI R1, 0x01} & \bub{00 LDI R1, 0x01} & \bub{00 LDI R1, 0x01} & 0 & 0 & -- & -- & -- \\
4 & \texttt{03 LDI R4, 0x20} & \texttt{02 LDI R3, 0x03} & \texttt{01 LDI R2, 0x10} & \texttt{00 LDI R1, 0x01} & \bub{00 LDI R1, 0x01} & 0 & 0 & -- & -- & -- \\
5 & \texttt{04 ST R3, [R2]} & \texttt{03 LDI R4, 0x20} & \texttt{02 LDI R3, 0x03} & \texttt{01 LDI R2, 0x10} & \texttt{00 LDI R1, 0x01} & 0 & 0 & -- & -- & R1$\leftarrow$01 \\
6 & \texttt{05 NOP} & \texttt{04 ST R3, [R2]} & \texttt{03 LDI R4, 0x20} & \texttt{02 LDI R3, 0x03} & \texttt{01 LDI R2, 0x10} & 0 & 0 & -- & -- & R2$\leftarrow$10 \\
7 & \texttt{06 ADD R1, R2} & \texttt{05 NOP} & \texttt{04 ST R3, [R2]} & \texttt{03 LDI R4, 0x20} & \texttt{02 LDI R3, 0x03} & 0 & 0 & MEM/WB & -- & R3$\leftarrow$03 \\
8 & \texttt{07 SUB R3, R1} & \texttt{06 ADD R1, R2} & \texttt{05 NOP} & \texttt{04 ST R3, [R2]} & \texttt{03 LDI R4, 0x20} & 0 & 0 & -- & -- & R4$\leftarrow$20; RAM[10]$\leftarrow$03 \\
9 & \texttt{08 LD R1, [R2]} & \texttt{07 SUB R3, R1} & \texttt{06 ADD R1, R2} & \texttt{05 NOP} & \texttt{04 ST R3, [R2]} & 0 & 0 & -- & -- & -- \\
10 & \texttt{09 ADD R3, R1} & \texttt{08 LD R1, [R2]} & \texttt{07 SUB R3, R1} & \texttt{06 ADD R1, R2} & \texttt{05 NOP} & 0 & 0 & -- & EX/MEM & -- \\
11 & \texttt{0A ADD R1, R2} & \texttt{09 ADD R3, R1} & \texttt{08 LD R1, [R2]} & \texttt{07 SUB R3, R1} & \texttt{06 ADD R1, R2} & \textbf{1} & 0 & -- & -- & R1$\leftarrow$11 \\
12 & \texttt{0A ADD R1, R2} & \texttt{09 ADD R3, R1} & \bub{09 ADD R3, R1} & \texttt{08 LD R1, [R2]} & \texttt{07 SUB R3, R1} & 0 & 0 & -- & -- & R3$\leftarrow$F2 \\
13 & \texttt{0B SUB R1, R3} & \texttt{0A ADD R1, R2} & \texttt{09 ADD R3, R1} & \bub{09 ADD R3, R1} & \texttt{08 LD R1, [R2]} & 0 & 0 & -- & MEM/WB & R1$\leftarrow$03 \\
14 & \texttt{0C ST R1, [R4]} & \texttt{0B SUB R1, R3} & \texttt{0A ADD R1, R2} & \texttt{09 ADD R3, R1} & \bub{09 ADD R3, R1} & 0 & 0 & -- & -- & -- \\
15 & \texttt{0D LD R5, [R4]} & \texttt{0C ST R1, [R4]} & \texttt{0B SUB R1, R3} & \texttt{0A ADD R1, R2} & \texttt{09 ADD R3, R1} & 0 & 0 & EX/MEM & MEM/WB & R3$\leftarrow$F5 \\
16 & \texttt{0E LDI R7, 0xF5} & \texttt{0D LD R5, [R4]} & \texttt{0C ST R1, [R4]} & \texttt{0B SUB R1, R3} & \texttt{0A ADD R1, R2} & 0 & 0 & EX/MEM & -- & R1$\leftarrow$13 \\
17 & \texttt{0F ST R5, [R7]} & \texttt{0E LDI R7, 0xF5} & \texttt{0D LD R5, [R4]} & \texttt{0C ST R1, [R4]} & \texttt{0B SUB R1, R3} & 0 & 0 & -- & -- & R1$\leftarrow$1E; RAM[20]$\leftarrow$1E \\
18 & \texttt{10 HALT} & \texttt{0F ST R5, [R7]} & \texttt{0E LDI R7, 0xF5} & \texttt{0D LD R5, [R4]} & \texttt{0C ST R1, [R4]} & 0 & 0 & -- & -- & -- \\
19 & \texttt{11 LDI R6, 0xFF} & \texttt{10 HALT} & \texttt{0F ST R5, [R7]} & \texttt{0E LDI R7, 0xF5} & \texttt{0D LD R5, [R4]} & 0 & 0 & MEM/WB & EX/MEM & R5$\leftarrow$1E \\
20 & \texttt{12 ST R6, [R7]} & \texttt{11 LDI R6, 0xFF} & \texttt{10 HALT} & \texttt{0F ST R5, [R7]} & \texttt{0E LDI R7, 0xF5} & 0 & \textbf{1} & -- & -- & R7$\leftarrow$F5; DISPLAY$\leftarrow$1E \\
21 & \textit{halted} & \bub{12 ST R6, [R7]} & \bub{11 LDI R6, 0xFF} & \texttt{10 HALT} & \texttt{0F ST R5, [R7]} & 0 & 0 & -- & -- & -- \\
22 & \textit{halted} & \bub{12 ST R6, [R7]} & \bub{12 ST R6, [R7]} & \bub{11 LDI R6, 0xFF} & \texttt{10 HALT} & 0 & 0 & -- & -- & -- \\
23 & \textit{halted} & \bub{12 ST R6, [R7]} & \bub{12 ST R6, [R7]} & \bub{12 ST R6, [R7]} & \bub{11 LDI R6, 0xFF} & 0 & 0 & -- & -- & -- \\
24 & \textit{halted} & \bub{12 ST R6, [R7]} & \bub{12 ST R6, [R7]} & \bub{12 ST R6, [R7]} & \bub{12 ST R6, [R7]} & 0 & 0 & -- & -- & -- \\

\end{tracetab}
\section{Full trace of \texttt{t2\_branch\_halt} (cycles 0--26)}
Wrong-path \texttt{HALT}s at \hex{0B} (squashed in cycle 13) and \hex{0F}/\hex{10} (squashed by \texttt{JMP} in cycle 17) are
ignored. The real \texttt{HALT} at \hex{13} is in EX in cycle 22 and kills the younger \texttt{ST R2,[R4]} at \hex{14}.
\begin{tracetab}
0 & \texttt{00 LDI R2, 0x21} & \bub{00 LDI R2, 0x21} & \bub{00 LDI R2, 0x21} & \bub{00 LDI R2, 0x21} & \bub{00 LDI R2, 0x21} & 0 & 0 & -- & -- & -- \\
1 & \texttt{00 LDI R2, 0x21} & \bub{00 LDI R2, 0x21} & \bub{00 LDI R2, 0x21} & \bub{00 LDI R2, 0x21} & \bub{00 LDI R2, 0x21} & 0 & 0 & -- & -- & -- \\
2 & \texttt{01 LDI R4, 0xF4} & \texttt{00 LDI R2, 0x21} & \bub{00 LDI R2, 0x21} & \bub{00 LDI R2, 0x21} & \bub{00 LDI R2, 0x21} & 0 & 0 & -- & -- & -- \\
3 & \texttt{02 LDI R3, 0x07} & \texttt{01 LDI R4, 0xF4} & \texttt{00 LDI R2, 0x21} & \bub{00 LDI R2, 0x21} & \bub{00 LDI R2, 0x21} & 0 & 0 & -- & -- & -- \\
4 & \texttt{03 LDI R1, 0x09} & \texttt{02 LDI R3, 0x07} & \texttt{01 LDI R4, 0xF4} & \texttt{00 LDI R2, 0x21} & \bub{00 LDI R2, 0x21} & 0 & 0 & -- & -- & -- \\
5 & \texttt{04 LDI R7, 0xF5} & \texttt{03 LDI R1, 0x09} & \texttt{02 LDI R3, 0x07} & \texttt{01 LDI R4, 0xF4} & \texttt{00 LDI R2, 0x21} & 0 & 0 & -- & -- & R2$\leftarrow$21 \\
6 & \texttt{05 SUB R1, R1} & \texttt{04 LDI R7, 0xF5} & \texttt{03 LDI R1, 0x09} & \texttt{02 LDI R3, 0x07} & \texttt{01 LDI R4, 0xF4} & 0 & 0 & -- & -- & R4$\leftarrow$F4 \\
7 & \texttt{06 BEQZ R1, target1} & \texttt{05 SUB R1, R1} & \texttt{04 LDI R7, 0xF5} & \texttt{03 LDI R1, 0x09} & \texttt{02 LDI R3, 0x07} & 0 & 0 & -- & -- & R3$\leftarrow$07 \\
8 & \texttt{07 ST R2, [R4]} & \texttt{06 BEQZ R1, target1} & \texttt{05 SUB R1, R1} & \texttt{04 LDI R7, 0xF5} & \texttt{03 LDI R1, 0x09} & 0 & 0 & MEM/WB & MEM/WB & R1$\leftarrow$09 \\
9 & \texttt{08 ADD R3, R4} & \texttt{07 ST R2, [R4]} & \texttt{06 BEQZ R1, target1} & \texttt{05 SUB R1, R1} & \texttt{04 LDI R7, 0xF5} & 0 & \textbf{1} & EX/MEM & -- & R7$\leftarrow$F5 \\
10 & \texttt{09 BNEZ R1, bad} & \bub{08 ADD R3, R4} & \bub{07 ST R2, [R4]} & \texttt{06 BEQZ R1, target1} & \texttt{05 SUB R1, R1} & 0 & 0 & -- & -- & R1$\leftarrow$00 \\
11 & \texttt{0A BNEZ R3, target2} & \texttt{09 BNEZ R1, bad} & \bub{08 ADD R3, R4} & \bub{07 ST R2, [R4]} & \texttt{06 BEQZ R1, target1} & 0 & 0 & -- & -- & -- \\
12 & \texttt{0B HALT} & \texttt{0A BNEZ R3, target2} & \texttt{09 BNEZ R1, bad} & \bub{08 ADD R3, R4} & \bub{07 ST R2, [R4]} & 0 & 0 & -- & -- & -- \\
13 & \texttt{0C ST R2, [R4]} & \texttt{0B HALT} & \texttt{0A BNEZ R3, target2} & \texttt{09 BNEZ R1, bad} & \bub{08 ADD R3, R4} & 0 & \textbf{1} & -- & -- & -- \\
14 & \texttt{0D ST R3, [R7]} & \bub{0C ST R2, [R4]} & \bub{0B HALT} & \texttt{0A BNEZ R3, target2} & \texttt{09 BNEZ R1, bad} & 0 & 0 & -- & -- & -- \\
15 & \texttt{0E JMP target3} & \texttt{0D ST R3, [R7]} & \bub{0C ST R2, [R4]} & \bub{0B HALT} & \texttt{0A BNEZ R3, target2} & 0 & 0 & -- & -- & -- \\
16 & \texttt{0F HALT} & \texttt{0E JMP target3} & \texttt{0D ST R3, [R7]} & \bub{0C ST R2, [R4]} & \bub{0B HALT} & 0 & 0 & -- & -- & -- \\
17 & \texttt{10 HALT} & \texttt{0F HALT} & \texttt{0E JMP target3} & \texttt{0D ST R3, [R7]} & \bub{0C ST R2, [R4]} & 0 & \textbf{1} & -- & -- & DISPLAY$\leftarrow$07 \\
18 & \texttt{11 LDI R6, 0xAA} & \bub{10 HALT} & \bub{0F HALT} & \texttt{0E JMP target3} & \texttt{0D ST R3, [R7]} & 0 & 0 & -- & -- & -- \\
19 & \texttt{12 ST R6, [R7]} & \texttt{11 LDI R6, 0xAA} & \bub{10 HALT} & \bub{0F HALT} & \texttt{0E JMP target3} & 0 & 0 & -- & -- & -- \\
20 & \texttt{13 HALT} & \texttt{12 ST R6, [R7]} & \texttt{11 LDI R6, 0xAA} & \bub{10 HALT} & \bub{0F HALT} & 0 & 0 & -- & -- & -- \\
21 & \texttt{14 ST R2, [R4]} & \texttt{13 HALT} & \texttt{12 ST R6, [R7]} & \texttt{11 LDI R6, 0xAA} & \bub{10 HALT} & 0 & 0 & EX/MEM & -- & -- \\
22 & \texttt{15 LDI R6, 0x55} & \texttt{14 ST R2, [R4]} & \texttt{13 HALT} & \texttt{12 ST R6, [R7]} & \texttt{11 LDI R6, 0xAA} & 0 & \textbf{1} & -- & -- & R6$\leftarrow$AA; DISPLAY$\leftarrow$AA \\
23 & \textit{halted} & \bub{15 LDI R6, 0x55} & \bub{14 ST R2, [R4]} & \texttt{13 HALT} & \texttt{12 ST R6, [R7]} & 0 & 0 & -- & -- & -- \\
24 & \textit{halted} & \bub{15 LDI R6, 0x55} & \bub{15 LDI R6, 0x55} & \bub{14 ST R2, [R4]} & \texttt{13 HALT} & 0 & 0 & -- & -- & -- \\
25 & \textit{halted} & \bub{15 LDI R6, 0x55} & \bub{15 LDI R6, 0x55} & \bub{15 LDI R6, 0x55} & \bub{14 ST R2, [R4]} & 0 & 0 & -- & -- & -- \\
26 & \textit{halted} & \bub{15 LDI R6, 0x55} & \bub{15 LDI R6, 0x55} & \bub{15 LDI R6, 0x55} & \bub{15 LDI R6, 0x55} & 0 & 0 & -- & -- & -- \\

\end{tracetab}
\end{landscape}
\section{Test program listings}
\textbf{t1\_hazards.asm}
\begin{listtab}
\texttt{00} & \texttt{1201} & I & \texttt{} & \texttt{LDI  R1, 0x01} & {\footnotesize } \\
\texttt{01} & \texttt{1410} & I & \texttt{} & \texttt{LDI  R2, 0x10} & {\footnotesize ALU operand and RAM address} \\
\texttt{02} & \texttt{1603} & I & \texttt{} & \texttt{LDI  R3, 0x03} & {\footnotesize } \\
\texttt{03} & \texttt{1820} & I & \texttt{} & \texttt{LDI  R4, 0x20} & {\footnotesize RAM address for 8.3} \\
\texttt{04} & \texttt{7680} & R & \texttt{} & \texttt{ST   R3, [R2]} & {\footnotesize RAM[0x10] = 0x03 (initialise before the load)} \\
\texttt{05} & \texttt{0000} & N & \texttt{} & \texttt{NOP} & {\footnotesize } \\
\texttt{06} & \texttt{3280} & R & \texttt{} & \texttt{ADD  R1, R2} & {\footnotesize R1 = 0x11} \\
\texttt{07} & \texttt{4640} & R & \texttt{} & \texttt{SUB  R3, R1} & {\footnotesize R3 = 0x03-0x11 = 0xF2   (R1 from EX/MEM)} \\
\texttt{08} & \texttt{6280} & R & \texttt{} & \texttt{LD   R1, [R2]} & {\footnotesize R1 = RAM[0x10] = 0x03} \\
\texttt{09} & \texttt{3640} & R & \texttt{} & \texttt{ADD  R3, R1} & {\footnotesize 1 bubble, R1 from MEM/WB: R3 = 0xF5} \\
\texttt{0A} & \texttt{3280} & R & \texttt{} & \texttt{ADD  R1, R2} & {\footnotesize R1 = 0x13} \\
\texttt{0B} & \texttt{42C0} & R & \texttt{} & \texttt{SUB  R1, R3} & {\footnotesize R1 = 0x13-0xF5 = 0x1E  (youngest writer of R1)} \\
\texttt{0C} & \texttt{7300} & R & \texttt{} & \texttt{ST   R1, [R4]} & {\footnotesize RAM[0x20] = 0x1E   (EX/MEM beats MEM/WB)} \\
\texttt{0D} & \texttt{6B00} & R & \texttt{} & \texttt{LD   R5, [R4]} & {\footnotesize R5 = 0x1E} \\
\texttt{0E} & \texttt{1EF5} & I & \texttt{} & \texttt{LDI  R7, 0xF5} & {\footnotesize } \\
\texttt{0F} & \texttt{7BC0} & R & \texttt{} & \texttt{ST   R5, [R7]} & {\footnotesize DISPLAY = 0x1E} \\
\texttt{10} & \texttt{F000} & N & \texttt{} & \texttt{HALT} & {\footnotesize } \\
\texttt{11} & \texttt{1CFF} & I & \texttt{} & \texttt{LDI  R6, 0xFF} & {\footnotesize younger than HALT: never executes} \\
\texttt{12} & \texttt{7DC0} & R & \texttt{} & \texttt{ST   R6, [R7]} & {\footnotesize } \\

\end{listtab}
\textbf{t2\_branch\_halt.asm}
\begin{listtab}
\texttt{00} & \texttt{1421} & I & \texttt{} & \texttt{LDI  R2, 0x21} & {\footnotesize LOCK|VALVE pattern (would open the valve)} \\
\texttt{01} & \texttt{18F4} & I & \texttt{} & \texttt{LDI  R4, 0xF4} & {\footnotesize Rvalve = actuator latch} \\
\texttt{02} & \texttt{1607} & I & \texttt{} & \texttt{LDI  R3, 0x07} & {\footnotesize } \\
\texttt{03} & \texttt{1209} & I & \texttt{} & \texttt{LDI  R1, 0x09} & {\footnotesize } \\
\texttt{04} & \texttt{1EF5} & I & \texttt{} & \texttt{LDI  R7, 0xF5} & {\footnotesize DISPLAY} \\
\texttt{05} & \texttt{4240} & R & \texttt{} & \texttt{SUB  R1, R1} & {\footnotesize R1 = 0} \\
\texttt{06} & \texttt{9209} & I & \texttt{} & \texttt{BEQZ R1, target1} & {\footnotesize taken, R1 forwarded from EX/MEM} \\
\texttt{07} & \texttt{7500} & R & \texttt{} & \texttt{ST   R2, [R4]} & {\footnotesize wrong path: must not write the valve} \\
\texttt{08} & \texttt{3700} & R & \texttt{} & \texttt{ADD  R3, R4} & {\footnotesize wrong path} \\
\texttt{09} & \texttt{A216} & I & \texttt{target1} & \texttt{BNEZ R1, bad} & {\footnotesize not taken (R1 = 0)} \\
\texttt{0A} & \texttt{A60D} & I & \texttt{} & \texttt{BNEZ R3, target2} & {\footnotesize taken (R3 = 7)} \\
\texttt{0B} & \texttt{F000} & N & \texttt{} & \texttt{HALT} & {\footnotesize wrong-path HALT: flushed, must not stop} \\
\texttt{0C} & \texttt{7500} & R & \texttt{} & \texttt{ST   R2, [R4]} & {\footnotesize wrong path} \\
\texttt{0D} & \texttt{77C0} & R & \texttt{target2} & \texttt{ST   R3, [R7]} & {\footnotesize DISPLAY = 0x07 (R3 untouched by wrong path)} \\
\texttt{0E} & \texttt{B011} & I & \texttt{} & \texttt{JMP  target3} & {\footnotesize } \\
\texttt{0F} & \texttt{F000} & N & \texttt{} & \texttt{HALT} & {\footnotesize wrong-path HALT behind a jump} \\
\texttt{10} & \texttt{F000} & N & \texttt{} & \texttt{HALT} & {\footnotesize } \\
\texttt{11} & \texttt{1CAA} & I & \texttt{target3} & \texttt{LDI  R6, 0xAA} & {\footnotesize } \\
\texttt{12} & \texttt{7DC0} & R & \texttt{} & \texttt{ST   R6, [R7]} & {\footnotesize DISPLAY = 0xAA: execution continued} \\
\texttt{13} & \texttt{F000} & N & \texttt{} & \texttt{HALT} & {\footnotesize real HALT} \\
\texttt{14} & \texttt{7500} & R & \texttt{} & \texttt{ST   R2, [R4]} & {\footnotesize younger than HALT: killed} \\
\texttt{15} & \texttt{1C55} & I & \texttt{} & \texttt{LDI  R6, 0x55} & {\footnotesize } \\
\texttt{16} & \texttt{1CBD} & I & \texttt{bad} & \texttt{LDI  R6, 0xBD} & {\footnotesize } \\
\texttt{17} & \texttt{7DC0} & R & \texttt{} & \texttt{ST   R6, [R7]} & {\footnotesize } \\
\texttt{18} & \texttt{F000} & N & \texttt{} & \texttt{HALT} & {\footnotesize } \\

\end{listtab}
\end{document}
